% Hand-written appendix.  Nothing generates this file.  It is \input after
% generated/appendix.tex, which has already issued \appendix, so the \section
% below is numbered as appendix material.  It carries the audit that used to
% live in counterexamples/lorentzian-jensen/paper.tex: the formulation the two
% statements were posed in, which steps of the proposed proof survive, the
% parity obstruction that kills the rest, and the enclosure procedure behind
% the witness of Theorem~\ref{thm:volume-distance} / Corollary~\ref{thm:lorentzian}.  Notation is
% spelled out rather than abbreviated: cxcase.sty binds \C to \mathbb C, so the
% cone is written \mathcal C, and \D, \Q, \T of the source file are written
% \mathrm D, \mathcal Q, \mathcal T here.
\section{Two new results for hyperbolic polynomials (author: S.~Sra)}
\label{sec:lorentzianaudit}
\noindent\dblue{\emph{These results are non-AI (by S.~Sra); including here to share with interested folks.}}

\vskip8pt

Replacing ``Lorentzian'' by ``hyperbolic'' turns the refuted principle into a
theorem.  Let $p:\R^n\to\R$ be homogeneous of degree $m\ge1$ and hyperbolic with
respect to $e$, normalized so that $p(e)>0$, and let $\Lambda_{++}$ be the open
hyperbolicity cone containing $e$.  By \citet{Garding1959},
$\Lambda_{++}$ is an open convex cone, $p$ is hyperbolic with respect to each of
its points, and $F:=p^{1/m}$ is positive, homogeneous of degree one and concave
on it.  The other external input is the Helton--Vinnikov theorem
\citep{HeltonVinnikov2007,LewisParriloRamana2005}, in the normalized form: a
ternary form $q$ of degree $m$, hyperbolic with respect to $\mathbf 1=(1,1,1)$
with $q(\mathbf 1)>0$, admits real symmetric $A_1,A_2,A_3$ with
$A_1+A_2+A_3=I_m$ and $q(r)=q(\mathbf 1)\det(r_1A_1+r_2A_2+r_3A_3)$.  Three
points at a time is exactly what a triangle inequality needs.

\begin{lemma}[Ternary restriction]\label{lem:lor-ternary}
Fix $x_1,x_2,x_3\in\Lambda_{++}$, put $T(r):=r_1x_1+r_2x_2+r_3x_3$ and
$q:=p\circ T$.  Then $q$ is homogeneous of degree $m$ and hyperbolic with
respect to $\mathbf 1$, and each coordinate vector $\varepsilon_i\in\R^3$ lies in
the hyperbolicity cone of $q$ containing $\mathbf 1$.
\end{lemma}
\begin{proof}
Homogeneity is immediate, and $w:=T(\mathbf 1)=x_1+x_2+x_3\in\Lambda_{++}$
because the cone is convex.  So $p$ is hyperbolic with respect to $w$, and for
every $r\in\R^3$ the univariate $\lambda\mapsto q(r+\lambda\mathbf 1)
=p(T(r)+\lambda w)$ has only real zeros; that is hyperbolicity with respect to
$\mathbf 1$.  For the second claim, the path
$\gamma_i(t)=(1-t)\varepsilon_i+t\mathbf 1$ satisfies
$T(\gamma_i(t))=x_i+t\sum_{j\ne i}x_j\in\Lambda_{++}$, so $q(\gamma_i(t))>0$ on
$[0,1]$: it joins $\varepsilon_i$ to $\mathbf 1$ inside the positive component.
\end{proof}

\begin{lemma}[Positive definite representatives]\label{lem:lor-pdrep}
In that setting the matrices $A_1,A_2,A_3$ supplied by Helton--Vinnikov are
positive definite, and for all $i,j$,
\[
p(x_i)=q(\mathbf 1)\det A_i,\qquad
p\Bigl(\frac{x_i+x_j}2\Bigr)=q(\mathbf 1)\det\Bigl(\frac{A_i+A_j}2\Bigr).
\]
\end{lemma}
\begin{proof}
Along the path of Lemma~\ref{lem:lor-ternary},
$L_i(t):=\sum_k\gamma_i(t)_kA_k$ satisfies
$q(\gamma_i(t))=q(\mathbf 1)\det L_i(t)>0$, so $L_i(t)$ is nonsingular on
$[0,1]$; the eigenvalues of a continuous path of nonsingular real symmetric
matrices cannot change sign, so the inertia is constant.  At $t=1$,
$L_i(1)=A_1+A_2+A_3=I_m\succ0$, hence $A_i=L_i(0)\succ0$.  The two identities
are the determinantal representation evaluated at $\varepsilon_i$ and at
$(\varepsilon_i+\varepsilon_j)/2$.
\end{proof}

\begin{theorem}[\statusproved: hyperbolic Jensen--Stein pseudometric]\label{thm:lor-hyperbolic}
Let $p$ be hyperbolic as above.  Then $\sqrt{\J_p}$, with
$\J_p(x,y)=\log p\bigl(\frac{x+y}2\bigr)-\frac12\log p(x)-\frac12\log p(y)$, is a
pseudometric on $\Lambda_{++}$.
\end{theorem}
\begin{proof}
Concavity of $F=p^{1/m}$ gives
$F(\frac{x+y}2)\ge\frac{F(x)+F(y)}2\ge\sqrt{F(x)F(y)}$, which after raising to
the $m$th power makes $\J_p\ge0$; symmetry and vanishing on the diagonal are
immediate.  For the triangle inequality fix $x_1,x_2,x_3\in\Lambda_{++}$ and take
$A_1,A_2,A_3\succ0$ from Lemma~\ref{lem:lor-pdrep}.  The factor $q(\mathbf 1)$
cancels, leaving
\[
\J_p(x_i,x_j)
=\log\frac{\det\bigl((A_i+A_j)/2\bigr)}{\sqrt{\det A_i\det A_j}}
=\delta_{\mathrm S}^2(A_i,A_j),
\]
the squared Stein divergence, whose square root is a metric on the positive
definite cone \citep{SraSdiv}.  Applying that triangle inequality to
$A_1,A_2,A_3$ gives the claim, and the triple was arbitrary.
\end{proof}

% The conclusion is a pseudometric rather than a metric because degenerate $p$
% collapse distances: for $p=\ell^m$ with $\ell$ linear, $\J_p(x,y)$ depends only on
% $\ell(x)$ and $\ell(y)$.  The same restriction principle transfers a determinant
% inequality to the cone.

\noindent Theorem~\ref{thm:lor-fourpoint} below proves a conjecture once shared with the author (S.~Sra) by~\citet{thomas2025}.
\begin{theorem}[\statusproved]\label{thm:lor-fourpoint}
With $F=p^{1/m}$, for all $a,b,c$ in the closed hyperbolicity cone,
\[
F(a+b)\,F(a+c)\;\ge\;F(a)\,F(a+b+c)+F(b)\,F(c).
\]
\end{theorem}
\begin{proof}
Write $\mathcal D(M)=\det(M)^{1/m}$.  We first prove the matrix form
$\mathcal D(A+B)\mathcal D(A+C)\ge\mathcal D(A)\mathcal D(A+B+C)+\mathcal D(B)\mathcal D(C)$
for $A,B,C\succeq0$.  For $A\succ0$, put $X=A^{-1/2}BA^{-1/2}$ and
$Y=A^{-1/2}CA^{-1/2}$; every term carries $\mathcal D(A)^2$, so it suffices to
show $\mathcal D(I+X)\mathcal D(I+Y)\ge\mathcal D(I+X+Y)+\mathcal D(X)\mathcal D(Y)$.
Order the eigenvalues of $X$ decreasingly and those of $Y$ increasingly.  Since a
product does not depend on the ordering,
$\mathcal D(I+X)\mathcal D(I+Y)=\bigl(\prod_i[(1+x_i+y_i)+x_iy_i]\bigr)^{1/m}$, and
Minkowski's determinant inequality
$\det(R+S)^{1/m}\ge\det R^{1/m}+\det S^{1/m}$ --- itself Jensen's inequality for
$t\mapsto\log(1+e^t)$ applied to the eigenvalues of $R^{-1/2}SR^{-1/2}$ ---
applied to the diagonal matrices $R=\diag(1+x_i+y_i)$ and $S=\diag(x_iy_i)$
gives
\[
\mathcal D(I+X)\mathcal D(I+Y)\ \ge\ \Bigl(\prod_i(1+x_i+y_i)\Bigr)^{1/m}
+\Bigl(\prod_ix_iy_i\Bigr)^{1/m}.
\]
Fiedler's upper bound \citep{Fiedler1971}, which pairs one decreasing spectrum
with the other increasing one, gives $\det(I+X+Y)\le\prod_i(1+x_i+y_i)$, while
the second term is $\mathcal D(X)\mathcal D(Y)$; a limit $\varepsilon\downarrow0$
in $A+\varepsilon I$ covers singular $A$.  For the cone statement, apply
Lemmas~\ref{lem:lor-ternary} and~\ref{lem:lor-pdrep} to $a,b,c\in\Lambda_{++}$ to
get $A,B,C\succ0$ with $p(ra+sb+tc)=q(\mathbf 1)\det(rA+sB+tC)$; then
$F(ra+sb+tc)=q(\mathbf 1)^{1/m}\mathcal D(rA+sB+tC)$ for $r,s,t\ge0$, and the
displayed inequality is the matrix one times $q(\mathbf 1)^{2/m}$.  Points of the
closed cone are handled by $u\mapsto u+\varepsilon e$, which lies in
$\Lambda_{++}$, and continuity.
\end{proof}

Nothing here collides with Theorem~\ref{thm:volume-distance} and
Corollary~\ref{thm:lorentzian}, and it is worth saying why not.  The witness
\eqref{eq:lorentzian-cubic} is strictly Lorentzian but \emph{not} hyperbolic:
$G(x,1)=\frac1{10}x^3+\frac32x^2+\frac{11}5x+1$ has discriminant
\[
18abcd-4b^3d+b^2c^2-4ac^3-27a^2d^2=-\frac{1499}{1250}<0
\qquad
\Bigl(a,b,c,d=\tfrac1{10},\tfrac32,\tfrac{11}5,1\Bigr),
\]
so it has one real and two complex conjugate roots.  A binary form is hyperbolic
with respect to some direction only if it splits into real linear factors --- each
factor $\ell_i$ contributes the root $-\ell_i(x)/\ell_i(e)$, which is real for
every real $x$ only when $\ell_i$ is real up to scale --- so $G$ is hyperbolic
with respect to no direction at all.  In two variables hyperbolicity is therefore
strictly stronger than strict Lorentzianity, and the witness lives in the gap.
That gap is the content of both results together: what makes the midpoint Jensen
gap a squared distance is the determinantal representation that hyperbolicity
buys through Helton--Vinnikov, and not log-concavity of the coefficients.

%%% Local Variables:
%%% mode: LaTeX
%%% TeX-master: "main"
%%% End:
